\documentclass[12pt]{ctexart} % 专门用于中文排版，基于 article 文档类。

%================== 页面布局 ==================%
\usepackage[a4paper, margin=2.5cm]{geometry} % A4纸，2.5cm页边距
\usepackage{filecontents} % 用于在tex文档内部创建文件

%================== 数学与符号 ==================%
\usepackage{amsmath} % 增强数学公式排版
\usepackage{amsthm} % 定理环境
\usepackage{mathrsfs} % 花体字体 \mathscr
\usepackage{amssymb} % 更多数学符号

%================== 图形与表格 ==================%
\usepackage{graphicx} % 插入图片
\usepackage{diagbox} % 斜线表头
\usepackage{booktabs} % 高质量表格线
\usepackage{makecell} % 增强表格单元格
\usepackage{array} % 增强表格和数组
\usepackage{multirow} % 合并多行
\usepackage{multicol} % 多列布局
\usepackage{longtable} % 跨页长表格
\usepackage{colortbl} % 表格上色
\usepackage{subcaption} % 子图/子表
\usepackage{wrapfig} % 图文环绕
\usepackage{overpic} % 图片上添加文本
\usepackage{float} % 浮动体控制，H选项
\usepackage{siunitx} % 数字和单位排版
\usepackage{tikz} % 矢量图形
\usetikzlibrary{shapes.geometric, arrows, positioning} % TikZ 库

%================== 代码高亮 ==================%
\usepackage{listings}
\lstset{
	numbers=left,
	numberstyle= \tiny,
	breaklines,
	extendedchars=false,
	basicstyle=\small\ttfamily,
	flexiblecolumns=true,
	showstringspaces=false,
	tabsize=4,
	frame=shadowbox,
	rulesepcolor= \color{ red!20!green!20!blue!20},
	xleftmargin=2em, xrightmargin=2em, aboveskip=1em,
	framexleftmargin=2em,
	escapeinside=``,
	keywordstyle=\color{blue!70},
	commentstyle=\color{red!50!green!50!blue!50},
}

\lstnewenvironment{Matlab}[1][]{
	\lstset{
		language=Matlab,
		deletekeywords={disp},
		keywordstyle=\color{blue},
		stringstyle=\color[rgb]{0.75,0,0.75},
		commentstyle=\songti\color[rgb]{0,0.5,0},
		#1
	}
}{}
\lstnewenvironment{Python}[1][]{
	\lstset{
		language=Python,
		keywordstyle=\color[RGB]{255,119,0},
		morekeywords={as},
		deletekeywords={print},
		keywordstyle=[2]\color[RGB]{144,0,144},
		morekeywords=[2]{print},
		stringstyle=\color[RGB]{0,170,0},
		commentstyle=\songti\color[RGB]{221,0,0},
		#1
	}
}{}

%================== 伪代码 ==================%
\usepackage[ruled,lined,vlined]{algorithm2e}
\renewcommand{\algorithmcfname}{算法}
\SetKwInput{KwIn}{输入}
\SetKwInput{KwOut}{输出}
\SetKwFor{For}{for}{do}{end for}
\SetKwIF{If}{ElseIf}{Else}{if}{then}{else if}{else}{end if}
\SetKw{KwWhile}{while}
\SetKw{KwDo}{do}
\SetKw{KwEndWhile}{end while}
\SetKw{Return}{return}
\SetCommentSty{rmfamily}
% 可跨页伪代码环境
\makeatletter
\newenvironment{breakablealgorithm}
{ \begin{center}
		\refstepcounter{algorithm}
		\hrule height.8pt depth0pt \kern2pt
		\renewcommand{\caption}[2][\relax]{%
			{\raggedright\textbf{\ALG@name~\thealgorithm} ##2\par}%
			\ifx\relax##1\relax
			\addcontentsline{loa}{algorithm}{\protect\numberline{\thealgorithm}##2}%
			\else
			\addcontentsline{loa}{algorithm}{\protect\numberline{\thealgorithm}##1}%
			\fi
			\kern2pt\hrule\kern2pt
		}
	}{
		\kern2pt\hrule\relax
	\end{center}
}
\makeatother

%================== 其它工具 ==================%
\usepackage{microtype}
\usepackage{setspace}
\usepackage{url}
\usepackage{emptypage}
\usepackage{enumerate}
\usepackage[section]{placeins}
\usepackage[dvipsnames, table]{xcolor}
\usepackage{tocloft}
\usepackage[colorlinks=true, linkcolor=black, citecolor=blue, urlcolor=blue, unicode=true]{hyperref}
\usepackage{cleveref}

%================== 参考文献设置 ==================%
\begin{filecontents}[overwrite]{aibib.bib}
	@misc{AITool-DeepSeek,
		note = {DeepSeek, DeepSeek-V3.1, 深度求索（DeepSeek）, 2025-09-07}
	}
\end{filecontents}

\usepackage[backend=biber,style=gb7714-2015]{biblatex}
\addbibresource{references.bib}
\addbibresource{aibib.bib}

%================== 全局设置 ==================%
% 【修改1】页码设置：页脚居中，阿拉伯数字，无页眉
\pagestyle{plain}

% 【修改2】正文默认字号设为小四号，行距为单倍行距
\zihao{-4}
\singlespacing

% 公式按章节编号
\numberwithin{equation}{section}

% 【修改3】章节标题格式（严格按照规范）
% 一级标题：四号黑体，居中，编号为"一、""二、"
% 二级标题：小四号黑体，左对齐
% 三级标题：小四号黑体，左对齐
\ctexset{
	section = {
		format = {\zihao{4}\heiti\centering},
		name = {,、},
		number = \chinese{section},
		beforeskip = {1.0ex plus 0.2ex minus .2ex},
		afterskip = {0.8ex plus 0.2ex minus .2ex},
	},
	subsection = {
		format = {\zihao{-4}\heiti\raggedright},
		beforeskip = {0.8ex plus 0.2ex minus .2ex},
		afterskip = {0.6ex plus 0.1ex minus .1ex},
	},
	subsubsection = {
		format = {\zihao{-4}\heiti\raggedright},
		beforeskip = {0.6ex plus 0.2ex minus .2ex},
		afterskip = {0.4ex plus 0.1ex minus .1ex},
	},
}

% 【修改4】摘要标题格式：四号黑体居中
\renewcommand{\abstractname}{\zihao{4}\heiti 摘要}

\date{}

% 自定义居中命令
\newcommand{\tableCenter}[1]{\makecell[c]{#1}}

%================== 正文 ==================%
% 【修改5】论文题目：三号黑体
\title{\heiti\zihao{3}标题}
\author{}

\begin{document}
	
	% 【修改6】第一页必须有页码，从1开始。删除 \thispagestyle{empty}
	\setcounter{page}{1}
	\date{\vspace{-2em}}
	\maketitle
	
	% 【修改7】调整标题与摘要之间的距离（更合理的间距）
	\vspace{-4em}
	% 摘要正文：小四号宋体
	\begin{abstract}{\zihao{-4}\songti} 
	
		随着
		
		\textbf{针对问题一：}
		
		\textbf{针对问题二：}
		
		\textbf{针对问题三：}
		
		本文特色在于：
		\\[\baselineskip]
		\noindent
		\textbf{关键词:}关键词1\qquad 关键词2\qquad 关键词3\qquad 关键词4
	\end{abstract}
	
	\newpage
	
	\section{问题重述}\label{sec:问题重述}
	
	\subsection{问题背景}\label{sec:问题背景}
	
	
	\subsection{问题提出}\label{sec:问题提出}
	
	\textbf{问题一：}
	
	
	\textbf{问题二：}
	
	
	\textbf{问题三：}
	
	
	\section{问题分析}\label{sec:问题分析}
	
	\subsection{问题一的分析}
	
	
	\subsection{问题二的分析}
	
	
	\subsection{问题三的分析}
	
	
	\section{模型假设}\label{sec:模型假设}
	
	\begin{enumerate}
		\item 假设1
		\item 假设2
	\end{enumerate}
	
	
	\section{符号说明}\label{sec:符号说明}
	
	% 【改进】符号说明表格：使用 tabular* 让表格宽度等于文本宽度，更美观
	\begin{table}[H]
		\centering
		\begin{tabular*}{\textwidth}{@{\extracolsep{\fill}} c l l}
			\toprule
			符号 & 含义 & 单位 \\
			\midrule
			$x$ & 变量含义 & \si{\metre} \\
			$v$ & 速度 & \si{\metre\per\second} \\
			\bottomrule
		\end{tabular*}
	\end{table}
	
	\section{数据预处理}\label{sec:数据预处理}
	
	
	\section{问题一的模型建立与求解}\label{sec:问题一的模型建立与求解}
	
	
	\section{问题二的模型建立与求解}\label{sec:问题二的模型建立与求解}
	
	% --- TikZ 流程图示例 ---
	\begin{figure}[H]
		\centering
		\begin{tikzpicture}[node distance=2cm, auto]
			\tikzstyle{startstop} = [rectangle, rounded corners, minimum width=3cm, minimum height=1cm, text centered, draw=black, fill=red!30]
			\tikzstyle{process} = [rectangle, minimum width=3cm, minimum height=1cm, text centered, draw=black, fill=orange!30]
			\tikzstyle{arrow} = [thick,->,>=stealth]
			
			\node (start) [startstop] {开始};
			\node (pro1) [process, below=of start] {处理步骤};
			\node (stop) [startstop, below=of pro1] {结束};
			
			\draw [arrow] (start) -- (pro1);
			\draw [arrow] (pro1) -- (stop);
		\end{tikzpicture}
		\caption{一个简单的 TikZ 流程图示例}
		\label{fig:tikz_example}
	\end{figure}
	
	
	% --- 多子图示例（取消注释使用）---
%	\begin{figure}[H]
%		\centering
%		\begin{minipage}{0.48\textwidth}
%			\centering
%			\includegraphics[width=\linewidth]{images/123.jpg}
%			\subcaption{子图标题1}
%			\label{fig:sub1}
%		\end{minipage}
%		\hfill
%		\begin{minipage}{0.48\textwidth}
%			\centering
%			\includegraphics[width=\linewidth]{images/123.jpg}
%			\subcaption{子图标题2}
%			\label{fig:sub2}
%		\end{minipage}
%		
%		\vspace{1em}
%		
%		\begin{minipage}{0.48\textwidth}
%			\centering
%			\includegraphics[width=\linewidth]{images/123.jpg}
%			\subcaption{子图标题3}
%			\label{fig:sub3}
%		\end{minipage}
%		\hfill
%		\begin{minipage}{0.48\textwidth}
%			\centering
%			\includegraphics[width=\linewidth]{images/123.jpg}
%			\subcaption{子图标题4}
%			\label{fig:sub4}
%		\end{minipage}
%		\caption{总标题}
%		\label{fig:main_figure}
%	\end{figure}
	
	
	% --- 单张图片示例（取消注释使用）---
%	\begin{figure}[H]
%		\centering
%		\includegraphics[width=0.7\textwidth]{images/123.jpg}
%		\caption{单张图片标题}
%		\label{fig:single_image}
%	\end{figure}
	
	
	% --- 伪代码示例 ---
	\begin{algorithm}[H]
		\caption{K-Means 聚类算法}\label{alg:kmeans}
		\KwIn{
			样本数据集 $D = \{x_1, x_2, \dots, x_m\}$\;
			簇的数量 $K$\;
			最大迭代次数 $N$\;
		}
		\KwOut{
			簇划分 $C = \{C_1, C_2, \dots, C_K\}$\;
		}
		\tcc{1. 初始化}
		从数据集 $D$ 中随机选择 $K$ 个样本作为初始簇中心 $\mu$\;
		初始化迭代次数 $n \leftarrow 0$\;
		\tcc{2. 开始迭代}
		\While{$n < N$ 且簇中心仍在变化}{
			\tcc{2.1 分配步骤（E-step）}
			\For{$k=1, \dots, K$}{
				初始化簇 $C_k = \emptyset$\;
			}
			\For{每个样本 $x_j \in D$}{
				计算 $x_j$ 与各簇中心的距离\;
				将 $x_j$ 划入最近的簇\;
			}
			\tcc{2.2 更新步骤（M-step）}
			\For{$k=1, \dots, K$}{
				更新簇中心 $\mu'_k = \frac{1}{|C_k|} \sum_{x \in C_k} x$\;
			}
			\tcc{2.3 判断是否收敛}
			\If{$\mu'_k = \mu_k$ for all $k$}{
				\tcc{簇中心不再变化，算法收敛}
				\textbf{break}\;
			}
			$\mu \leftarrow \mu'$\;
			$n \leftarrow n + 1$\;
		}
		\Return 最终簇划分 $C$\;
	\end{algorithm}
	
	
	\section{问题三的模型建立与求解}\label{sec:问题三的模型建立与求解}
	
	
	\section{模型的评价与推广}
	
	\subsection{模型的优点}
	
	\begin{enumerate}
		\item 优点1
		\item 优点2
	\end{enumerate}
	
	
	\subsection{模型的缺点}
	
	\begin{enumerate}
		\item 缺点1
		\item 缺点2
	\end{enumerate}
	
	
	\subsection{模型的推广}
	
	
	%================== 参考文献 ==================%
	\nocite{AITool-DeepSeek}
	\printbibliography
	
	%================== 附录 ==================%
	\newpage
	\appendix
	
	% 【修改8】附录标题格式调整：取消居中，改为左对齐四号黑体（与正文章节区分）
	\ctexset{section={
			format={\zihao{4}\heiti\raggedright}
	}}
	
	\begin{center}
		\heiti\zihao{3} 附\hspace{1pc}录
	\end{center}
	
	\section{支撑材料文件列表}\label{sec:支撑材料文件列表}
	
	\begin{table}[H]
		\centering
		\label{tab:支撑材料文件列表}
		\renewcommand{\arraystretch}{1.2}
		\begin{tabular*}{\textwidth}{l l @{\extracolsep{\fill}} l}
			\toprule
			\textbf{文件名} & \textbf{类型} & \textbf{简介} \\
			\midrule
			第一问代码.py & Python 代码文件 & 问题一模型建立与求解代码 \\
			第二问代码.m  & Matlab 代码文件 & 问题二模型建立与求解代码 \\
			第三问代码.py & Python 代码文件 & 问题三模型建立与求解代码 \\
			人工智能工具使用说明.pdf & PDF 文件 & 人工智能工具的使用说明 \\
			数据预处理.xlsx & Excel 文件 & 原始数据的预处理结果 \\
			\bottomrule
		\end{tabular*}
	\end{table}
	
	\section{第一问代码（Python）}\label{sec:第一问代码}

	
	\section{第二问代码（Matlab）}\label{sec:第二问代码}
	\begin{Python}

	\end{Python}
	\section{第三问代码（Python）}\label{sec:第三问代码}
	\begin{Python}
# 这是一个 Python 代码示例
def hello_world():
    print("Hello, World!")

if __name__ == "__main__":
    hello_world()
	\end{Python}
	
\end{document}